\documentclass[12pt]{article}
\usepackage[utf8]{inputenc}
\usepackage[T1]{fontenc}
\usepackage{hyperref}
\title{Tonsawang Language Resources for K-12 Curriculum Development}
\author{Compiled by Claude Sonnet 5 (claude-sonnet-5) for Andr\'as Kornai}
\date{}
\begin{document}
\maketitle

\section{Introduction}
Tonsawang (ISO 639-3 tnw), also called Tombatu or Toundanow, is an endangered Austronesian language of the Minahasan branch, spoken in Southeast Minahasa, North Sulawesi, Indonesia. Speaker estimates vary: Omniglot gives about 9,000, and Wikipedia cites 20,000 from a 1981 source. Speakers recorded for documentation were mainly older adults. Younger people are shifting to Manado Malay. The community's own name for the language is Toundanow. Brickell describes it as the most divergent of the Minahasan languages. It is written in the Latin alphabet, with the glottal stop as an apostrophe.

\section{Primary Educational Resources}
The search summary mentions revitalization efforts, including teaching in schools, speech contests and use in churches. I did not see a primary source that gives details, and no primers or textbooks turned up.

\section{Linguistic and Reference Materials}
\begin{itemize}
\item Wikipedia, ``Tonsawang language'': \url{https://www.wikipedia.com/wiki/Tonsawang}.
\item Omniglot, Tonsawang: \url{https://omniglot.com/writing/tonsawang.htm}.
\item Journal article on Tonsawang: \url{https://journal.unj.ac.id/unj/index.php/isllae/article/download/12023/7397/26190}. I did not confirm its subject from the snippet.
\item Language Documentation and Description article: \url{https://lddjournal.org/article/1203/galley/2447/download}. I did not confirm its subject from the snippet.
\end{itemize}

\section{Religious and Community Resources}
Church use is mentioned in the search summary, but no specific scripture or community organization turned up.

\section{Digital and Audio Resources}
\begin{itemize}
\item ``Tonsawang: a collaborative multimedia project documenting an endangered language of North Sulawesi'', an ELDP-funded archive of narratives, songs and proverbs recorded 2016--2018: \url{https://eldp.access.preservica.com/dk0434/}.
\item Endangered Languages Project entry (sample): \url{https://endangeredlanguages.com/lang/3524/samples/11427?hl=de}.
\end{itemize}

\section{Curriculum Development Context}
The ELDP corpus is the strongest existing resource and could supply texts and audio for materials. Children are not acquiring the language, so any curriculum would be revitalization-oriented, built on the archive, with community involvement. Orthography needs to be settled. No existing school materials were found.

\end{document}